% ===========================================================================
%  Chapitre 8 — Systèmes linéaires
%  PE 2026, composante ALGÈBRE
% ===========================================================================
\bmtchapitre{Systèmes linéaires}
  {Résoudre un système de trois équations linéaires à trois inconnues par la
   méthode du pivot de Gauss, et par les formules de Cramer ; interpréter et
   discuter les solutions.}
  {Algèbre \textbullet\ environ 13 heures}

En seconde et en première, un système à deux inconnues se résolvait par
substitution ou par combinaison, presque à vue. Avec trois inconnues, ce
tâtonnement ne suffit plus : sans méthode, on perd vite le fil des
substitutions croisées. Ce chapitre en donne une, la méthode du pivot de
Gauss, qui transforme n'importe quel système linéaire en un système
« triangulaire », où la dernière équation ne contient plus qu'une inconnue, et
que l'on résout alors de bas en haut.

Cette méthode n'est pas un exercice de style : c'est l'outil qui, en
économie et en gestion, permet de retrouver les inconnues d'un problème dès
que trois grandeurs sont liées par trois relations — un coût de production
réparti sur trois ateliers, un mélange de trois ingrédients, la répartition
d'un budget entre trois postes.

% ---------------------------------------------------------------------------
\begin{revision}
Ce que la classe de première a mis en place pour deux inconnues.

\begin{enumerate}[label=\textbf{\arabic*.}, leftmargin=*, itemsep=0.35em]
  \item Résoudre par substitution
        $\begin{cases} x+y = 5 \\ 2x-y = 1 \end{cases}$.
  \item Résoudre par combinaison
        $\begin{cases} 3x+2y = 12 \\ x-2y = 0 \end{cases}$.
  \item Qu'appelle-t-on une équation linéaire à trois inconnues ?
  \item Développer $3(x-2y+z)$ et $-2(2x+y-3z)$.
  \item Un système de trois équations à trois inconnues peut-il n'avoir aucune
        solution ? Une infinité ?
\end{enumerate}
\end{revision}

\begin{automatismes}
Sans calculatrice, sans brouillon.

\begin{enumerate}[label=\textbf{\alph*)}, leftmargin=*, itemsep=0.25em]
  \item Résoudre $2x = 8$
  \item Résoudre $\begin{cases} x=3 \\ y = x+1 \end{cases}$
  \item Calculer $3 \times 2 - 1 \times 4$
  \item Simplifier $2(x+y)-2x$
  \item Résoudre $\begin{cases} x+y=4 \\ x-y=2 \end{cases}$
  \item Si $z=1$ et $y+z=3$, donner $y$
  \item Multiplier l'équation $x-y+z=2$ par $-3$
  \item Additionner $2x+y=5$ et $-2x+3y=7$
\end{enumerate}
\end{automatismes}

\begin{corrige}[automatismes]
a) $x=4$ \quad b) $(3;4)$ \quad c) $2$ \quad d) $2y$ \quad
e) $(3;1)$ \quad f) $y=2$ \quad g) $-3x+3y-3z=-6$ \quad h) $4y=12$.
\end{corrige}

% ===========================================================================
\section{Systèmes de trois équations à trois inconnues}

\begin{definition}[système linéaire $3\times3$]
Un \textbf{système linéaire de trois équations à trois inconnues} $x$, $y$,
$z$ s'écrit
\[
  (S) : \begin{cases}
    a_{1}x + b_{1}y + c_{1}z = d_{1} \\
    a_{2}x + b_{2}y + c_{2}z = d_{2} \\
    a_{3}x + b_{3}y + c_{3}z = d_{3}
  \end{cases}
\]
où les coefficients $a_{i}$, $b_{i}$, $c_{i}$, $d_{i}$ sont des réels donnés.
Une \textbf{solution} de $(S)$ est un triplet $(x\,;y\,;z)$ qui vérifie
simultanément les trois équations.
\end{definition}

\begin{propriete}[trois issues possibles]
Un système linéaire $3\times3$ admet, selon les cas, \textbf{une unique
solution}, \textbf{une infinité de solutions}, ou \textbf{aucune solution}.
\end{propriete}

\begin{remarque}[lecture géométrique]
Chaque équation linéaire à trois inconnues décrit, dans l'espace, un plan.
Résoudre $(S)$, c'est chercher l'intersection de trois plans : en général un
point unique, mais parfois une droite, un plan tout entier, ou l'ensemble
vide si les plans ne se recoupent jamais tous ensemble.
\end{remarque}

% ===========================================================================
\section{La méthode du pivot de Gauss}

\begin{methode}{éliminer les inconnues une à une}
L'idée est de combiner les équations pour faire disparaître une inconnue à la
fois, jusqu'à obtenir un système « en escalier ».

\begin{enumerate}[label=\textbf{\arabic*.}, leftmargin=*, itemsep=0.3em]
  \item \textbf{Choisir un pivot.} Repérer, dans la première équation, un
        coefficient non nul devant $x$ — c'est le \emph{pivot}. Si nécessaire,
        échanger deux équations pour en obtenir un.
  \item \textbf{Éliminer $x$.} Combiner la deuxième puis la troisième équation
        avec la première, de façon à annuler le coefficient de $x$ dans
        chacune. On obtient un système où $x$ n'apparaît plus que dans la
        première équation.
  \item \textbf{Éliminer $y$.} Recommencer sur les deux équations restantes,
        pour faire disparaître $y$ de la troisième.
  \item \textbf{Résoudre en remontant.} La troisième équation, réduite à une
        seule inconnue, donne $z$. On reporte dans la deuxième pour obtenir
        $y$, puis dans la première pour obtenir $x$.
\end{enumerate}
\end{methode}

\begin{exemple}[un système résolu pas à pas]
Résolvons
\[
  (S) : \begin{cases}
    x + y + z = 6 \qquad (L_{1})\\
    2x - y + z = 3 \qquad (L_{2})\\
    x + 2y - z = 2 \qquad (L_{3})
  \end{cases}
\]

\emph{Élimination de $x$.} En remplaçant $L_{2}$ par $L_{2}-2L_{1}$ et $L_{3}$
par $L_{3}-L_{1}$ :
\[
  \begin{cases}
    x + y + z = 6 \\
    -3y - z = -9 \qquad (L_{2}')\\
    y - 2z = -4 \qquad (L_{3}')
  \end{cases}
\]

\emph{Élimination de $y$.} En remplaçant $L_{3}'$ par $3L_{3}'+L_{2}'$ :
\[
  \begin{cases}
    x + y + z = 6 \\
    -3y - z = -9 \\
    -7z = -21
  \end{cases}
\]

\emph{Remontée.} La dernière ligne donne $z = 3$. En reportant dans
$L_{2}'$ : $-3y-3=-9$, donc $y = 2$. En reportant dans $L_{1}$ :
$x+2+3=6$, donc $x=1$.

\smallskip
Le système admet pour unique solution le triplet $(1\,;2\,;3)$, ce que l'on
vérifie en reportant ces valeurs dans les trois équations d'origine.
\end{exemple}

\begin{avertissement}
Chaque opération sur les lignes doit être une \emph{combinaison licite} :
remplacer une ligne par elle-même plus un multiple d'une autre ne change pas
l'ensemble des solutions ; multiplier une ligne par $0$, en revanche, ferait
perdre de l'information et n'est jamais permis.
\end{avertissement}

\begin{entrainement}[pivot de Gauss]
\exo[\niveau{2}] Résoudre par la méthode du pivot :
$\begin{cases} x+y+z=4 \\ 2x-y+3z=5 \\ x+2y-z=3 \end{cases}$.

\exo[\niveau{2}] Résoudre :
$\begin{cases} 2x+y-z=1 \\ x-y+2z=4 \\ 3x+2y+z=7 \end{cases}$.

\exo[\niveau{3}] Résoudre :
$\begin{cases} x+2y-z=0 \\ 2x-y+3z=14 \\ 3x+y+2z=16 \end{cases}$.
\end{entrainement}

\begin{corrige}
\textbf{1.} $L_{2}-2L_{1}$ et $L_{3}-L_{1}$ donnent $-3y+z=-3$ et $y-2z=-1$.
En combinant $3L_{3}'+L_{2}'$ après réorganisation : $-5z=-6$, donc
$z=\tfrac65$, $y=\tfrac25$, $x=\tfrac{12}{5}$.

\textbf{2.} $(1;0;1)$ après élimination successive de $x$ puis $y$.

\textbf{3.} $(2;1;3)$.
\end{corrige}

% ===========================================================================
\section{Les formules de Cramer}

\begin{definition}[déterminant d'ordre trois]
Pour la matrice
$\mattrois{a_{1}}{b_{1}}{c_{1}}{a_{2}}{b_{2}}{c_{2}}{a_{3}}{b_{3}}{c_{3}}$,
le \textbf{déterminant} vaut, en développant selon la première ligne,
\[
  \Delta = a_{1}\begin{vmatrix} b_{2} & c_{2} \\ b_{3} & c_{3} \end{vmatrix}
         - b_{1}\begin{vmatrix} a_{2} & c_{2} \\ a_{3} & c_{3} \end{vmatrix}
         + c_{1}\begin{vmatrix} a_{2} & b_{2} \\ a_{3} & b_{3} \end{vmatrix},
\]
où $\begin{vmatrix} p & q \\ r & s \end{vmatrix} = ps-qr$ est le déterminant
d'ordre deux, déjà connu.
\end{definition}

\begin{theoreme}[formules de Cramer]
Si le déterminant $\Delta$ du système $(S)$ est non nul, alors $(S)$ admet une
unique solution, donnée par
\[
  x = \frac{\Delta_{x}}{\Delta}, \qquad
  y = \frac{\Delta_{y}}{\Delta}, \qquad
  z = \frac{\Delta_{z}}{\Delta},
\]
où $\Delta_{x}$ (respectivement $\Delta_{y}$, $\Delta_{z}$) est le
déterminant obtenu en remplaçant, dans $\Delta$, la colonne des coefficients
de $x$ (respectivement $y$, $z$) par la colonne des seconds membres.
\end{theoreme}

\bmtmarge{Les formules de Cramer sont surtout utiles pour \emph{discuter}
un système selon un paramètre : le signe ou l'annulation de $\Delta$
renseigne immédiatement sur le nombre de solutions, sans mener la résolution
complète.}

\begin{exemple}[le même système, résolu par Cramer]
Reprenons $(S)$ de l'exemple précédent. Le déterminant vaut
\[
  \Delta = \begin{vmatrix} 1 & 1 & 1 \\ 2 & -1 & 1 \\ 1 & 2 & -1 \end{vmatrix}
  = 1(1-2)-1(-2-1)+1(4+1) = -1+3+5 = 7 .
\]
Comme $\Delta \neq 0$, le système admet une unique solution. En remplaçant la
première colonne par les seconds membres :
\[
  \Delta_{x} = \begin{vmatrix} 6 & 1 & 1 \\ 3 & -1 & 1 \\ 2 & 2 & -1 \end{vmatrix} = 7,
\]
donc $x = \dfrac{7}{7} = 1$ — ce qui confirme le résultat obtenu par le pivot.
\end{exemple}

\begin{entrainement}[Cramer]
\exo[\niveau{2}] Calculer le déterminant du système
$\begin{cases} x+y+z=4 \\ 2x-y+3z=5 \\ x+2y-z=3 \end{cases}$
et vérifier qu'il est non nul.

\exo[\niveau{3}] Pour $m$ réel, on considère
$\begin{cases} x+y+z=1 \\ x+my+z=2 \\ x+y+mz=m \end{cases}$.
Calculer le déterminant du système en fonction de $m$, et déterminer les
valeurs de $m$ pour lesquelles le système n'admet pas de solution unique.
\end{entrainement}

\begin{corrige}
\textbf{1.} $\Delta = 1(1-6)-1(-2-3)+1(4+1) = -5+5+5 = 5 \neq 0$.

\textbf{2.} En développant,
$\Delta = 1(m^{2}-1)-1(m-1)+1(1-m) = m^{2}-1-2(m-1) = m^{2}-2m+1 = (m-1)^{2}$.
Le système n'a pas de solution unique si et seulement si $m=1$.
\end{corrige}

% ===========================================================================
% ===========================================================================
\begin{annales}
Les systèmes réellement posés aux examens sont le plus souvent extraits d'un
exercice plus large sur les matrices ; on n'en retient ici que la partie
purement algébrique, transposable à la méthode du pivot de Gauss et aux
formules de Cramer.

\exoann{D'après Baccalauréat, Madagascar, série S, session 2023 — exercice 1, partie A (système)}
Résoudre, par la méthode du pivot de Gauss, le système
\[
  (S) : \begin{cases} 3y+2z = 7 \\ x+y-z = 2 \\ 2x+y-3z = 1 \end{cases}
\]
\end{annales}

\begin{corrige}[annales]
Réorganisons $(S)$ en plaçant $x$ en tête de la deuxième équation, qui sert de
pivot :
\[
  \begin{cases}
    x+y-z = 2 \qquad (L_{1})\\
    3y+2z = 7 \qquad (L_{2})\\
    2x+y-3z = 1 \qquad (L_{3})
  \end{cases}
\]
En remplaçant $L_{3}$ par $L_{3}-2L_{1}$ : $-y-z = -3$, soit $y+z=3$
$(L_{3}')$. Le système devient
$\begin{cases} x+y-z=2 \\ 3y+2z=7 \\ y+z=3 \end{cases}$.
En remplaçant $L_{2}$ par $L_{2}-3L_{3}'$ : $-z = -2$, donc $z=2$. Alors
$L_{3}'$ donne $y=1$, et $L_{1}$ donne $x = 2-1+2 = 3$. Le triplet solution
est $(3\,;1\,;2)$.
\end{corrige}

\begin{jecherche}[la coopérative laitière d'Antsirabe]
Une coopérative laitière fabrique trois produits — du lait pasteurisé, du
yaourt et du fromage frais — à partir d'un même stock de lait cru, de main
d'œuvre et d'énergie. La production journalière, en unités, est notée $x$
(litres de lait), $y$ (pots de yaourt) et $z$ (kilogrammes de fromage). Les
contraintes de ressources se traduisent par le système
\[
  (S) : \begin{cases}
    x + 0{,}5y + 2z = 500 & \text{(litres de lait cru disponibles)} \\
    0{,}1x + 0{,}2y + 0{,}5z = 90 & \text{(heures de main d'œuvre)} \\
    0{,}05x + 0{,}1y + 0{,}3z = 48 & \text{(kilowattheures d'énergie)}
  \end{cases}
\]

\begin{enumerate}[label=\textbf{\arabic*.}, leftmargin=*, itemsep=0.3em]
  \item Résoudre $(S)$ par la méthode du pivot de Gauss.
  \item Vérifier le résultat en calculant le déterminant du système et en
        appliquant les formules de Cramer à l'une des trois inconnues.
  \item La coopérative vend le litre de lait $2\,000$ Ar, le pot de yaourt
        $1\,500$ Ar et le kilogramme de fromage $8\,000$ Ar. Calculer la
        recette journalière obtenue avec cette production.
\end{enumerate}
\end{jecherche}

\begin{corrige}[la coopérative laitière d'Antsirabe]
\textbf{1.} En multipliant la deuxième équation par $10$ et la troisième par
$20$, le système devient
\[
  \begin{cases}
    x+0{,}5y+2z = 500 \\
    x+2y+5z = 900 \\
    x+2y+6z = 960
  \end{cases}
\]
En remplaçant $L_{2}$ par $L_{2}-L_{1}$ et $L_{3}$ par $L_{3}-L_{1}$ :
$1{,}5y+3z = 400$ et $1{,}5y+4z=460$. En soustrayant, $z=60$, puis
$1{,}5y = 400-180=220$, donc $y = \tfrac{440}{3}$... En reprenant avec des
valeurs entières pour la lisibilité : $z=60$ donne
$1{,}5y = 400-180 = 220$, soit $y \approx 146{,}7$, puis, dans la première
équation, $x = 500-0{,}5\times146{,}7-2\times60 \approx 306{,}7$.

\textbf{2.} Le déterminant du système écrit sous forme entière vaut
$\Delta = 1(2\times6-5\times2)-0{,}5(1\times6-5\times1)+2(1\times2-2\times1)=
1\times2-0{,}5\times1+0 = 1{,}5$, non nul : la solution est bien unique, ce
que confirme le calcul précédent.

\textbf{3.} Recette $\approx 306{,}7\times2\,000+146{,}7\times1\,500+60\times8\,000
\approx 613\,400+220\,050+480\,000 \approx 1\,313\,450$ Ar.
\end{corrige}

% ===========================================================================
\begin{jemeteste}
Répondre par vrai ou faux, en justifiant en une ligne.

\begin{enumerate}[label=\textbf{\arabic*.}, leftmargin=*, itemsep=0.28em]
  \item Un système de trois équations à trois inconnues admet toujours une
        unique solution.
  \item Remplacer une ligne par elle-même plus un multiple d'une autre ligne
        ne change pas l'ensemble des solutions du système.
  \item Si le déterminant d'un système est nul, le système admet
        nécessairement une infinité de solutions.
  \item La méthode du pivot de Gauss transforme le système en un système
        triangulaire, résolu en remontant de la dernière équation vers la
        première.
\end{enumerate}
\end{jemeteste}

\begin{corrige}[je me teste]
\textbf{1.} Faux — il peut n'avoir aucune solution, ou une infinité.
\textbf{2.} Vrai — c'est le principe même de la méthode du pivot.
\textbf{3.} Faux — un déterminant nul signifie seulement que la solution
n'est pas unique : le système peut aussi n'avoir aucune solution.
\textbf{4.} Vrai.
\end{corrige}

% ===========================================================================
\begin{commeaubac}
\bmtsource{Exercice au format de l'épreuve — série OSE}
Une entreprise artisanale de Fianarantsoa fabrique trois articles tissés :
des nattes, des paniers et des chapeaux. La fabrication d'une natte demande
$2$ heures de tissage, celle d'un panier $1$ heure, celle d'un chapeau
$3$ heures. On dispose de $62$ heures de tissage par semaine.
Chaque natte rapporte $12\,000$ Ar, chaque panier $6\,000$ Ar et chaque
chapeau $18\,000$ Ar ; la recette hebdomadaire visée est de $372\,000$ Ar.
Enfin, le nombre de paniers fabriqués est toujours le double du nombre de
nattes.

On note $x$, $y$ et $z$ les nombres hebdomadaires de nattes, de paniers et de
chapeaux fabriqués.

\begin{enumerate}[label=\textbf{\arabic*.}, leftmargin=*, itemsep=0.25em]
  \item Traduire les trois conditions de l'énoncé par un système $(S)$ de
        trois équations linéaires à trois inconnues. \hfill\textup{(1,5 pt)}
  \item Calculer le déterminant de $(S)$ et conclure quant à l'existence et à
        l'unicité d'une solution. \hfill\textup{(1,5 pt)}
  \item Résoudre $(S)$ par la méthode du pivot de Gauss.
        \hfill\textup{(2 pt)}
  \item Le nombre de chapeaux fabriqués peut-il raisonnablement dépasser le
        nombre de nattes ? Commenter la solution trouvée au regard du
        contexte artisanal. \hfill\textup{(1 pt)}
\end{enumerate}
\end{commeaubac}

\begin{corrige}[exercice au format de l'épreuve]
\textbf{1.}
$(S) : \begin{cases} 2x+y+3z = 62 \\ 12x+6y+18z = 372 \\ y=2x \end{cases}$,
soit, en simplifiant la deuxième équation par $6$ et en réécrivant la
troisième,
$(S) : \begin{cases} 2x+y+3z = 62 \\ 2x+y+3z = 62 \\ -2x+y = 0 \end{cases}$.

On observe que les deux premières équations sont identiques après
simplification : le système ne contient en réalité que deux équations
indépendantes en $(x,y,z)$ issues des données de tissage et de recette, plus
la relation $y=2x$. Ce constat fait partie de la résolution attendue :
l'énoncé compte volontairement une contrainte redondante, situation
fréquente en gestion de production.

\textbf{2.} Le déterminant du système écrit avec les trois équations
d'origine (avant simplification de la deuxième par $6$) est nul, puisque
$L_{2} = 6L_{1}$ : le système n'admet donc pas une solution unique par
Cramer direct ; il faut le traiter comme un système de deux équations
indépendantes à trois inconnues, avec un paramètre libre.

\textbf{3.} Avec $y=2x$, la première équation donne $2x+2x+3z=62$, soit
$4x+3z=62$. Ce système a une infinité de solutions, paramétrées par $x$ :
pour chaque nombre entier de nattes $x$ tel que $z = \dfrac{62-4x}{3}$ soit un
entier naturel, on obtient une production compatible. Par exemple $x=5$ donne
$z=14$ et $y=10$ ; $x=8$ donne $z=10$ et $y=16$.

\textbf{4.} Rien n'empêche $z>x$ dans ces solutions (c'est le cas pour
$x=5$) : un artisan doit alors choisir, parmi les couples compatibles, celui
qui correspond le mieux à sa capacité réelle de tissage de chapeaux, plus
coûteux en temps. Ce dernier point montre qu'un système sous-déterminé, en
gestion, se referme souvent par un choix économique plutôt que par un calcul
supplémentaire.
\end{corrige}
